\section{Limitations}
\textbf{Data and evaluation.} The benchmark consists of 53 patients and 1015 criteria; the fixed split leaves only eight test patients. The 95\,\% bootstrap intervals in Table~\ref{tab:main} are wide (e.g.\ accuracy 0.60--0.80), several classes have a handful of test examples (\emph{not included}: 7, inclusion N/A: 3) and the test split contains \emph{no} \emph{excluded} example, so the detection of excluded patients is not evaluated. The patients are short curated case descriptions (at most 14 sentences and 372 tokens), not full electronic health records, so no conclusion about long, messy notes can be drawn; for the same reason truncation never occurred, and the behaviour for notes beyond 555 tokens is untested.

\textbf{Modelling.} Only two training configurations were tried and each was trained once with a single seed, so differences of a few points between runs (e.g.\ 0.728 versus 0.733 accuracy) are within noise; the macro-F1 gain of the balanced run (0.467 to 0.534 on validation) rests on the exclusion classes, which have few validation examples. Other open models (Llama-3.2-1B, Gemma-2-2B) were rejected on resource and licensing grounds without being trained, so we make no claim that Qwen2.5-1.5B is the best choice, only that it fits the budget. The model returns a label and sentence IDs; the displayed explanation is a template that quotes the cited sentences and is not a free-text rationale, so it cannot explain \emph{why} the sentences support the label. The model is weakest exactly where a screening tool needs care: it labels most missing-information exclusion criteria \emph{not excluded} (3 of 28 NEI rows correct) and omits evidence in 42 of 58 wrong predictions. The trial-level outcome in the application is a transparent rule-based aggregation that we did not validate against any ground truth.

\textbf{Inference experiments.} All timings come from one laptop GPU, one software stack and small samples (12 task prompts, 4 long prompts, three repeats for the KV benchmark); they are indicative, not definitive. Fp16 decoding is not bit-reproducible between the cached and uncached paths, and the observed divergences are explained by exact logit ties (Section~\ref{sec:kvres}); an fp32 test was possible only for the 0.5\,B model. The tiled attention is a forward-pass educational implementation (no backward pass, no fused kernel) and is slower than the dense reference (Section~\ref{sec:attres}); we do not claim a speed benefit, only exactness and lower memory. Speculative decoding is implemented for greedy decoding only; its slowdown is specific to this hardware, an eager decoding loop and a draft that is only 19\,\% cheaper than the target, and part of the gap between predicted and measured speed-up is not isolated. The application was tested through scripted and API calls and screenshots, not with a formal usability study.

\textbf{Intended use.} The system is an educational prototype. It must not be used for enrolment or any clinical decision.

\section{Conclusion}
We built and evaluated a complete, locally executable pipeline for criterion-level clinical-trial eligibility matching on a 6\,GB GPU. LoRA fine-tuning of Qwen2.5-1.5B-Instruct on only 37 patients raised held-out accuracy from 0.058 to 0.718 and macro-F1 from 0.084 to 0.604 while returning cited sentence IDs (evidence micro-F1 0.531 overall, 0.636 on rows with expert evidence), and class-balanced oversampling was the decisive design choice for the exclusion classes. The honest picture has two sides: the model beats a majority baseline overall and on inclusion criteria, but on exclusion criteria it does not beat it in accuracy and it largely fails to say ``not enough information'', treating silence in the note as eligibility. Our hand-written decoding components behaved as theory predicts: the explicit KV cache is exact in fp32 and 2.2--7$\times$ faster in fp16 (with divergences traced to exact logit ties), tiled online-softmax attention matches a float64 reference to $4\times10^{-5}$ while cutting attention memory by up to 158$\times$ at the price of speed in plain PyTorch, and cache-synchronised greedy speculative decoding reproduces greedy output exactly with 62--95\,\% draft acceptance but is 0.7--0.8$\times$ as fast on this hardware because the draft is almost as costly per step as the target. Future work includes a larger and more balanced evaluation (cross-validation over patients), calibration of the NEI decision, a cheaper or compiled draft, and generating rationales with verifiable citations.
